\section{Discussion}
\label{sec:discussion}

As shown in Cross-Track Synthesis, GPT-4.1-mini demonstrates substantially superior generation quality compared to Claude-3.5-Haiku across all experimental tracks. The paradoxical relationship between training utility and detectability proves instructive: GPT's superior training effectiveness stems from closer distributional alignment to real spam, meaning GPT-generated content preserves the lexical patterns, feature co-occurrences, and statistical regularities that characterize genuine spam emails. This distributional fidelity enables classifiers to learn decision boundaries that generalize effectively to real spam detection tasks. However, this same proximity makes GPT synthetic spam more detectable classifiers trained exclusively on real spam can identify GPT-generated content because it shares sufficient statistical overlap with the training distribution.

In contrast, Claude generates synthetic spam with greater distributional divergence from real spam, resulting in both reduced training utility and paradoxically lower detectability. This suggests that high-quality synthetic data generation requires achieving distributional fidelity sufficient for effective classifier training while maintaining enough divergence to evade detection—a delicate tension. GPT navigates this balance more successfully through subtle linguistic realism rather than crude distributional mimicry, generating spam that "sounds natural" while preserving the statistical patterns classifiers need for effective learning. For attackers seeking training augmentation, GPT provides superior utility; for attackers prioritizing evasion, both models offer pathways but through different mechanisms.

As demonstrated in Cross-Track Synthesis, Random Forest consistently outperforms SVM across all experimental tracks. This architectural advantage stems from ensemble methods' inherent resistance to distributional shifts between real and synthetic spam. Random Forest's bootstrap aggregation mechanism trains multiple decision trees on random subsets of features and samples, enabling the ensemble to capture diverse statistical patterns while avoiding overfitting to specific distributional artifacts that may characterize synthetic data. When encountering synthetic spam during training or cross-model scenarios during testing, individual trees may respond differently to distributional anomalies, but the voting mechanism averages out these sensitivities, producing stable aggregate predictions.

For operational deployment, this architectural difference carries concrete implications. Organizations prioritizing stability and cross-model robustness should deploy Random Forest architectures, accepting slightly higher computational costs in exchange for consistent performance across varied synthetic augmentation strategies. Conversely, organizations requiring predictable degradation patterns for capacity planning may prefer SVM despite higher sensitivity, as the linear degradation enables reliable performance forecasting and intervention thresholds.

Synthetic augmentation from one LLM improves detection of another LLM's spam, demonstrating the potential for "fighting fire with fire" defensive strategies. However, the counterintuitive finding that weak prompt training data provides larger improvements than strong prompt data reveals that diversity rather than quality drives cross-model generalization. This pattern suggests that lower-quality, more divergent synthetic spam better captures the full feature space of AI-generated content across varied generation strategies, enabling robust detection of synthetic spam regardless of the specific LLM or prompt used for generation.

This finding fundamentally inverts the quality-diversity trade-off compared to training augmentation scenarios. For offensive training augmentation (Track 1), attackers should prioritize high-fidelity synthetic data matching real spam distributions. For defensive cross-model augmentation (Track 2b), defenders should prioritize diverse, lower-quality synthetic data capturing broader distributional variations. This asymmetry creates actionable guidance: defenders seeking maximum robustness against unknown LLM-generated spam should intentionally generate diverse, varied synthetic training data rather than optimizing for realism, accepting lower individual sample quality in exchange for broader feature space coverage.

The precision-recall asymmetry observed in Track 1 reveals an additional limitation of synthetic training data. Classifiers trained on pure synthetic spam maintain high precision while suffering severe recall degradation, becoming overly conservative—correctly identifying spam when detected but missing many actual spam instances. This pattern suggests synthetic spam captures surface-level lexical features enabling spam recognition but fails to represent the full distributional diversity of real spam, causing classifiers to reject borderline cases. Organizations must account for this bias when establishing acceptance criteria, potentially requiring hybrid training strategies mixing real and synthetic data to preserve recall performance rather than relying exclusively on synthetic augmentation.

These findings illuminate the fundamental asymmetry in Generative AI's impact on spam detection. For attackers, LLMs provide substantial advantages: as shown in Cross-Track Synthesis, GPT-4.1-mini generates synthetic spam maintaining high training effectiveness when substituted for real spam, while remaining moderately detectable by traditional defenses. The low generation cost, minimal technical barriers, and demonstrated evasion capabilities establish LLMs as accessible, effective offensive tools. Attackers can leverage superior models and optimal prompts to maximize either training augmentation or evasion, depending on strategic objectives.

For defenders, benefits prove more constrained: synthetic augmentation improves detection but requires cross-model training data, architectural choices favoring robustness over performance, and acceptance of precision-recall trade-offs. The substantial quality gap between different LLM models suggests attackers selecting optimal models gain advantages defenders cannot fully mitigate through augmentation alone. Defenders' primary advantage lies in ensemble methods' inherent cross-model robustness, providing baseline resilience independent of specific augmentation strategies. Overall, Generative AI appears to favor attackers in spam generation while providing defenders with limited but measurable countermeasures, shifting the traditional offense-defense balance toward attackers unless defenders proactively adopt ensemble architectures and cross-model augmentation strategies.

% \subsection{Limitations and Future Directions}

Several limitations constrain generalizability and suggest directions for future research. Our TF-IDF feature representation captures lexical patterns but may miss semantic characteristics where deep learning models with contextualized embeddings (BERT, RoBERTa) might exhibit different sensitivity profiles to synthetic spam. The CEAS-08 dataset represents 2008 spam characteristics; evolving adversarial tactics and linguistic patterns may interact differently with modern LLM-generated synthetic data, warranting evaluation on contemporary spam corpora. Our fixed 10\% spam ratio represents one point on the class imbalance spectrum; more severe imbalance or near-balanced scenarios may produce different degradation patterns and optimal augmentation strategies.
The two LLM models we evaluate (GPT-4.1-mini, Claude-3.5-Haiku) represent mid-tier commercial APIs optimized for cost-performance trade-offs. Several unexplored model categories may alter the offense-defense balance: (1) \textit{open-source LLMs} such as Llama 3, Mistral, and Qwen offering lower costs, local deployment capabilities, and greater privacy control; (2) \textit{broader price-performance tiers} ranging from budget models (GPT-3.5-turbo, Claude-3-Haiku) to frontier models (GPT-4o, Claude-3.5-Sonnet, o1) with potentially different quality-cost trade-offs for both attackers and defenders; and (3) \textit{specialized generation approaches} including traditional methods (template-based, Markov chains), domain-specific fine-tuned models, or hybrid approaches combining multiple generation strategies. These alternatives may provide superior cost-efficiency, controllability, or evasion capabilities depending on attacker objectives and defender constraints.

Track 2 experiments focus on cross-model detection scenarios, while same-model detection (attackers and defenders using identical LLMs) remains unexplored and may exhibit different patterns due to shared distributional characteristics. Finally, our spam detection context may not generalize to other cybersecurity domains where synthetic data characteristics and detection requirements differ fundamentally from email classification tasks. Future work should explore these dimensions to establish comprehensive understanding of LLM-generated synthetic data across diverse security contexts.
